\documentclass[11pt,a4paper]{article}

\usepackage[a4paper,margin=24mm]{geometry}
\usepackage{fontspec}
\usepackage[spanish]{babel}
\usepackage{microtype}
\usepackage{xcolor}
\usepackage{colortbl}
\usepackage{setspace}
\usepackage{titlesec}
\usepackage{enumitem}
\usepackage{booktabs}
\usepackage{tabularx}
\usepackage{array}
\usepackage{hyperref}
\usepackage{tcolorbox}
\usepackage{fancyhdr}
\usepackage{lastpage}
\usepackage{ragged2e}
\usepackage{xparse}

\tcbuselibrary{skins,breakable,raster}

% Tipografias libres y faciles de conseguir en distribuciones TeX
\setmainfont{TeX Gyre Pagella}
\setsansfont{TeX Gyre Heros}
\setmonofont{Latin Modern Mono}

% Paleta inspirada en Claude / Anthropic
\definecolor{claudeParchment}{HTML}{F5F4ED}
\definecolor{claudeIvory}{HTML}{FAF9F5}
\definecolor{claudeTerracotta}{HTML}{C96442}
\definecolor{claudeCoral}{HTML}{D97757}
\definecolor{claudeInk}{HTML}{141413}
\definecolor{claudeCharcoal}{HTML}{30302E}
\definecolor{claudeWarmText}{HTML}{4D4C48}
\definecolor{claudeOliveGray}{HTML}{5E5D59}
\definecolor{claudeStone}{HTML}{87867F}
\definecolor{claudeBorder}{HTML}{E8E6DC}
\definecolor{claudeBorderSoft}{HTML}{F0EEE6}
\definecolor{claudeWarmSilver}{HTML}{B0AEA5}
\definecolor{claudeForest}{HTML}{5A6E63}
\definecolor{claudeError}{HTML}{B53333}

\pagecolor{claudeParchment}
\color{claudeOliveGray}

\hypersetup{
  colorlinks=true,
  linkcolor=claudeTerracotta,
  urlcolor=claudeTerracotta,
  citecolor=claudeTerracotta,
  pdfborder={0 0 0}
}

\setlength{\parindent}{0pt}
\setlength{\parskip}{0.78em}
\setstretch{1.28}
\setlist[itemize]{leftmargin=1.4em,itemsep=0.35em,topsep=0.3em}
\setlist[enumerate]{leftmargin=1.6em,itemsep=0.35em,topsep=0.3em}
\setlength{\headheight}{18pt}
\setlength{\tabcolsep}{6pt}
\emergencystretch=2em

\newcolumntype{Y}{>{\RaggedRight\arraybackslash}X}

% Metadatos editables
\newcommand{\DocKicker}{Editorial Template}
\newcommand{\DocTitle}{Plantilla editorial inspirada en Claude}
\newcommand{\DocSubtitle}{Sistema reusable para informes, propuestas, memorandos y documentos largos con tono cálido y sereno}
\newcommand{\DocAuthor}{Tu Nombre}
\newcommand{\DocRole}{Dirección de diseño / redacción}
\newcommand{\DocDate}{4 de mayo de 2026}
\newcommand{\DocRef}{Ref. CLAUDE-LATEX-002}
\newcommand{\DocFooter}{Claude Editorial LaTeX Template}

\newcommand{\BodyCopy}{\color{claudeOliveGray}}
\newcommand{\DisplayFont}{\rmfamily}
\newcommand{\UIFont}{\sffamily}

% Estilo de secciones
\titleformat{\section}
  {\DisplayFont\Large\bfseries\color{claudeInk}}
  {}
  {0pt}
  {}
  [\vspace{0.38em}{\color{claudeBorder}\titlerule}]

\titleformat{\subsection}
  {\DisplayFont\large\bfseries\color{claudeWarmText}}
  {}
  {0pt}
  {}

\titlespacing*{\section}{0pt}{1.8em}{0.8em}
\titlespacing*{\subsection}{0pt}{1.1em}{0.45em}

% Cabecera y pie
\pagestyle{fancy}
\fancyhf{}
\renewcommand{\headrulewidth}{0pt}
\renewcommand{\footrulewidth}{0pt}
\fancyhead[L]{\UIFont\footnotesize\color{claudeWarmText}\DocKicker}
\fancyhead[R]{\UIFont\footnotesize\color{claudeStone}\DocRef}
\fancyfoot[L]{\UIFont\footnotesize\color{claudeStone}\DocFooter}
\fancyfoot[R]{\UIFont\footnotesize\color{claudeTerracotta}\thepage/\pageref{LastPage}}

% Base para cajas
\tcbset{
  enhanced,
  boxrule=0.6pt,
  colback=claudeIvory,
  colframe=claudeBorder,
  arc=3mm,
  left=10pt,
  right=10pt,
  top=8pt,
  bottom=8pt
}

\newtcolorbox{claudehero}{
  breakable,
  colback=claudeIvory,
  colframe=claudeBorderSoft,
  boxrule=0.8pt,
  arc=4mm,
  left=14pt,
  right=14pt,
  top=14pt,
  bottom=14pt,
  borderline west={2pt}{0pt}{claudeTerracotta},
  coltext=claudeWarmText
}

\newtcolorbox{claudenote}{
  breakable,
  colback=claudeIvory,
  colframe=claudeBorder,
  borderline west={1.6pt}{0pt}{claudeTerracotta},
  coltext=claudeWarmText
}

\newtcolorbox{claudequote}{
  breakable,
  colback=claudeIvory,
  colframe=claudeBorderSoft,
  borderline west={1.6pt}{0pt}{claudeCoral},
  coltext=claudeInk
}

\newtcolorbox{claudecard}{
  breakable,
  colback=claudeIvory,
  colframe=claudeBorder,
  coltext=claudeWarmText
}

\newtcolorbox{claudeaside}[1]{
  breakable,
  colback=claudeIvory,
  colframe=claudeBorder,
  title=\UIFont\bfseries\color{claudeInk}#1,
  coltitle=claudeInk,
  coltext=claudeWarmText
}

\newtcolorbox{claudecallout}[1]{
  breakable,
  colback=claudeIvory,
  colframe=claudeBorder,
  title=\UIFont\bfseries\color{claudeInk}#1,
  coltitle=claudeInk,
  coltext=claudeInk
}

\newtcolorbox{claudedecision}[1]{
  breakable,
  colback=claudeIvory,
  colframe=claudeTerracotta,
  boxrule=0.8pt,
  title=\UIFont\bfseries\color{claudeInk}#1,
  coltitle=claudeInk,
  coltext=claudeWarmText,
  borderline west={2pt}{0pt}{claudeTerracotta}
}

\newtcolorbox{claudechecklist}[1]{
  breakable,
  colback=claudeIvory,
  colframe=claudeBorderSoft,
  title=\UIFont\bfseries\color{claudeInk}#1,
  coltitle=claudeInk,
  coltext=claudeWarmText,
  borderline west={2pt}{0pt}{claudeForest}
}

\newtcolorbox{claudedarkpanel}[1]{
  breakable,
  colback=claudeInk,
  colframe=claudeCharcoal,
  boxrule=0.6pt,
  arc=4mm,
  title=\UIFont\bfseries\color{claudeIvory}#1,
  coltitle=claudeIvory,
  coltext=claudeWarmSilver,
  left=14pt,
  right=14pt,
  top=14pt,
  bottom=14pt
}

\newcommand{\ClaudeRule}{\par\noindent{\color{claudeBorder}\rule{\textwidth}{0.6pt}}\par}
\newcommand{\ClaudeSubtleRule}{\par\noindent{\color{claudeBorderSoft}\rule{\textwidth}{0.4pt}}\par}

\NewDocumentCommand{\ClaudeBadge}{O{claudeTerracotta} m}{%
  \tcbox[
    on line,
    colback=claudeIvory,
    colframe=claudeBorderSoft,
    boxrule=0.5pt,
    arc=2.5mm,
    left=6pt,
    right=6pt,
    top=3pt,
    bottom=3pt
  ]{\UIFont\scriptsize\color{#1}#2}%
}

\NewDocumentCommand{\ClaudeLead}{m}{%
  {\fontsize{13.4}{20}\selectfont\color{claudeWarmText}#1\par}%
}

\NewDocumentCommand{\ClaudeStat}{m m m}{%
  \begin{tcolorbox}[
    colback=claudeIvory,
    colframe=claudeBorderSoft,
    boxrule=0.6pt,
    arc=3mm,
    left=10pt,
    right=10pt,
    top=10pt,
    bottom=10pt
  ]
  {\UIFont\scriptsize\color{claudeStone}#1\par}
  \vspace{0.25em}
  {\DisplayFont\fontsize{23}{26}\selectfont\color{claudeInk}#2\par}
  \vspace{0.25em}
  {\small\color{claudeWarmText}#3\par}
  \end{tcolorbox}%
}

\NewDocumentCommand{\ClaudeFact}{m m}{%
  {\UIFont\scriptsize\color{claudeStone}#1\par}
  {\DisplayFont\normalsize\color{claudeInk}#2\par}
  \vspace{0.7em}
}

\NewDocumentCommand{\ClaudeTimelineItem}{m m m}{%
  \begin{tabularx}{\textwidth}{@{}p{0.2\textwidth}@{\hspace{1em}}Y@{}}
    {\UIFont\small\color{claudeStone}#1} &
    {\DisplayFont\normalsize\bfseries\color{claudeInk}#2\par
    {\small\color{claudeWarmText}#3}}
  \end{tabularx}
}

\NewDocumentCommand{\ClaudeSignature}{m m}{%
  \vspace{1.5em}
  {\color{claudeBorder}\rule{48mm}{0.5pt}\par}
  \vspace{0.7em}
  {\DisplayFont\large\color{claudeInk}#1\par}
  {\UIFont\small\color{claudeStone}#2\par}
}

\begin{document}
\BodyCopy

{\small\UIFont\color{claudeTerracotta}\DocKicker}

{\fontsize{31}{35}\selectfont\DisplayFont\color{claudeInk}\DocTitle\par}

{\large\UIFont\color{claudeWarmText}\DocSubtitle\par}

\vspace{0.65em}

{\UIFont\small\color{claudeStone}\DocAuthor\hspace{1em}·\hspace{1em}\DocRole\hspace{1em}·\hspace{1em}\DocDate\par}

\vspace{0.8em}
\ClaudeBadge{Claude-inspired}
\hspace{0.35em}\ClaudeBadge[claudeForest]{Warm editorial}
\hspace{0.35em}\ClaudeBadge[claudeWarmText]{Reusable system}

\vspace{0.55em}
\ClaudeRule
\vspace{0.45em}

\begin{claudehero}
\ClaudeLead{Esta plantilla ya no funciona como una simple página de ejemplo. Está planteada como un sistema editorial reusable con suficientes piezas para componer documentos reales: informes, memorandos, propuestas, notas estratégicas, resúmenes ejecutivos y textos largos con un tono más humano y más deliberado.}

{\UIFont\small\color{claudeStone}La referencia visual principal es \texttt{design-systems/claude/DESIGN.md}, reforzada con el criterio de \texttt{design-systems/warm-editorial/DESIGN.md}.}
\end{claudehero}

\vspace{0.85em}

\begin{tcbraster}[
  raster columns=3,
  raster equal height=rows,
  raster column skip=8pt,
  raster before skip=0pt,
  raster after skip=0pt
]
\ClaudeStat{Jerarquía}{Serif + sans}{Titulares con cadencia editorial y metadatos en una sans funcional y discreta.}
\ClaudeStat{Color}{Terracotta + warm neutrals}{Un acento medido sobre pergamino cálido y grises con subtono orgánico.}
\ClaudeStat{Componentes}{Más piezas disponibles}{Portada, notas, decisiones, checklist, comparativas, cronología, bloque oscuro y cierre.}
\end{tcbraster}

\section{Qué incluye esta plantilla}

\ClaudeLead{La idea no es copiar una interfaz, sino trasladar a documento el lenguaje visual de Claude: calma, lectura pausada, serif con gravedad, bordes suaves y contraste sin frialdad tecnológica.}

\begin{minipage}[t]{0.62\textwidth}
Esta versión está pensada para que reutilices más de una capa de diseño al mismo tiempo: la portada,
los resúmenes, los módulos laterales, las tablas comparativas, las métricas y los cortes de atmósfera.
No tienes por qué usarlo todo en cada documento; la gracia está en que ya existe una pequeña biblioteca
de piezas y no una única maqueta cerrada.

\begin{claudenote}
\UIFont\small
\textbf{\color{claudeInk}Resumen ejecutivo.}
Si la primera versión era una buena prueba de tono, esta ya sirve como base de trabajo real:
puedes mezclar módulos sin que el documento pierda consistencia.
\end{claudenote}
\end{minipage}\hfill
\begin{minipage}[t]{0.34\textwidth}
\begin{claudeaside}{Tokens rápidos}
\ClaudeFact{Fondo}{Parchment \#F5F4ED}
\ClaudeFact{Superficie}{Ivory \#FAF9F5}
\ClaudeFact{Acento}{Terracotta \#C96442}
\ClaudeFact{Texto}{Ink \#141413}
\ClaudeFact{Secundario}{Olive Gray \#5E5D59}
\ClaudeFact{Borde}{Warm Border \#E8E6DC}
\end{claudeaside}
\end{minipage}

\section{Piezas disponibles}

\begin{claudecallout}{Inventario de componentes}
\ClaudeBadge{Portada editorial}
\hspace{0.3em}\ClaudeBadge[claudeForest]{Resumen ejecutivo}
\hspace{0.3em}\ClaudeBadge[claudeWarmText]{Métricas}
\hspace{0.3em}\ClaudeBadge{Nota}
\hspace{0.3em}\ClaudeBadge[claudeTerracotta]{Decisión}
\hspace{0.3em}\ClaudeBadge[claudeForest]{Checklist}
\hspace{0.3em}\ClaudeBadge[claudeWarmText]{Tabla comparativa}
\hspace{0.3em}\ClaudeBadge{Línea temporal}
\hspace{0.3em}\ClaudeBadge[claudeForest]{Bloque oscuro}
\hspace{0.3em}\ClaudeBadge[claudeWarmText]{Firma / cierre}

\vspace{0.8em}
\begin{itemize}
\item informes breves con presencia visual;
\item propuestas narrativas con datos puntuales;
\item documentación institucional menos rígida;
\item textos largos con cambios de ritmo entre secciones.
\end{itemize}
\end{claudecallout}

\clearpage

\section{Bloques editoriales}

\ClaudeLead{Aquí es donde la plantilla se vuelve útil de verdad: no solo hay estilo, hay piezas con función clara.}

\begin{tcbraster}[
  raster columns=2,
  raster equal height=rows,
  raster column skip=10pt,
  raster before skip=0pt,
  raster after skip=0pt
]
\begin{claudenote}
\textbf{\sffamily Nota editorial.}

Úsala para añadir contexto, antecedentes, aclaraciones o marco conceptual.
\end{claudenote}

\begin{claudequote}
\textbf{\sffamily Cita o idea central.}

\itshape
La sensación general debe ser la de una página bien editada: cálida, serena, precisa y con suficiente aire para que el contenido respire.
\end{claudequote}

\begin{claudedecision}{Decisión / recomendación}
Este bloque sirve para fijar postura, aprobación, criterio o conclusión. Tiene más peso visual y por eso conviene reservarlo para decisiones de verdad.
\end{claudedecision}

\begin{claudechecklist}{Checklist de uso}
\begin{itemize}
\item dejar respirar el papel;
\item usar el terracota con moderación;
\item mantener la serif para los momentos de autoridad;
\item usar la sans solo donde aporta claridad funcional.
\end{itemize}
\end{claudechecklist}
\end{tcbraster}

\section{Comparativas y evidencia}

\begin{claudecallout}{Matriz de composición}
\begingroup
\arrayrulecolor{claudeBorder}
\begin{tabularx}{\textwidth}{@{}>{\DisplayFont\color{claudeInk}}lYYY@{}}
\toprule
{\color{claudeWarmText}Pieza} & {\color{claudeWarmText}Función} & {\color{claudeWarmText}Tono} & {\color{claudeWarmText}Uso recomendado} \\
\midrule
Portada & Abrir el documento & Editorial y sereno & Informes, propuestas, notas con entidad propia \\
Nota & Dar contexto breve & Calmado y auxiliar & Avisos, marco legal, matices, antecedentes \\
Decisión & Fijar postura & Firme y visible & Acuerdos, directrices, recomendaciones \\
Bloque oscuro & Cambiar la atmósfera & Más denso y dramático & Capítulos, manifiestos, resúmenes de alto peso \\
Cronología & Ordenar hitos & Claro y narrativo & Planes de fase, evolución, fases de proyecto \\
\bottomrule
\end{tabularx}
\endgroup
\end{claudecallout}

\section{Línea temporal}

\begin{claudecard}
\ClaudeTimelineItem{Semana 1}{Definición visual}{Ajustar la portada, la jerarquía tipográfica y el contraste general del documento.}
\vspace{0.8em}
\ClaudeSubtleRule
\vspace{0.8em}
\ClaudeTimelineItem{Semana 2}{Sistema de componentes}{Añadir cajas reutilizables, métricas, tablas, bloques laterales y módulos de decisión.}
\vspace{0.8em}
\ClaudeSubtleRule
\vspace{0.8em}
\ClaudeTimelineItem{Semana 3}{Composición avanzada}{Introducir patrones de dos columnas, cortes oscuros y páginas con más ritmo editorial.}
\vspace{0.8em}
\ClaudeSubtleRule
\vspace{0.8em}
\ClaudeTimelineItem{Semana 4}{Producción real}{Aplicar la plantilla a documentos concretos y consolidar variantes según caso de uso.}
\end{claudecard}

\clearpage

\begin{claudedarkpanel}{Sección oscura opcional}
{\DisplayFont\fontsize{24}{28}\selectfont\color{claudeIvory}Un documento largo también necesita cambios de atmósfera.\par}

\vspace{0.5em}
{\color{claudeWarmSilver}Este módulo sirve para aperturas de capítulo, manifiestos breves, resúmenes de alto nivel o páginas donde interese concentrar más peso visual sin abandonar el sistema Claude.}

\vspace{0.9em}
\ClaudeBadge[claudeWarmSilver]{Dark section}
\hspace{0.35em}\ClaudeBadge[claudeWarmSilver]{Chapter opener}
\hspace{0.35em}\ClaudeBadge[claudeWarmSilver]{High-contrast pause}
\end{claudedarkpanel}

\section{Patrones de composición}

\ClaudeLead{Además de las piezas sueltas, ya hay suficientes bloques como para repetir patrones de página según el tipo de documento.}

\begin{tcbraster}[
  raster columns=3,
  raster equal height=rows,
  raster column skip=8pt,
  raster before skip=0pt,
  raster after skip=0pt
]
\begin{claudeaside}{Apertura}
Titular, subtítulo, metadatos, badges y resumen ejecutivo en una primera página que respira.
\end{claudeaside}

\begin{claudeaside}{Desarrollo}
Secciones con una idea principal, una cita, una tabla puntual o un módulo lateral para no caer en monotonía.
\end{claudeaside}

\begin{claudeaside}{Cierre}
Decisión, próximos pasos, firma y, si conviene, una última llamada de tono institucional.
\end{claudeaside}
\end{tcbraster}

\section{Cierre}

Si vas a reutilizar esta plantilla, no hace falta que conserves el ejemplo completo. Puedes tomar solo las piezas que te interesen: portada + nota para un memo breve, o portada + comparativa + cronología + decisión para un informe más estructurado.

\begin{claudedecision}{Próximo paso recomendado}
El siguiente paso natural sería separar esto en dos archivos:
\texttt{claude-editorial-base.tex} para la estructura reusable y
\texttt{claude-editorial-demo.tex} para este catálogo visual de componentes.
\end{claudedecision}

\ClaudeSignature{\DocAuthor}{\DocRole}

\end{document}
